% Generated draft: requires mathematical and source review.
\documentclass[11pt]{article}
\usepackage[margin=25mm]{geometry}
\usepackage{fontspec}
\setmainfont{Latin Modern Roman}
\usepackage{amsmath,amssymb,mathtools,unicode-math}
\setmathfont{Latin Modern Math}
\usepackage{enumitem,needspace,xurl,hyperref}
\setlength{\parindent}{0pt}
\setlength{\parskip}{5pt}
\setlength{\emergencystretch}{3em}
\newcommand{\N}{\mathbb{N}}
\newcommand{\Z}{\mathbb{Z}}
\newcommand{\Q}{\mathbb{Q}}
\newcommand{\R}{\mathbb{R}}
\newcommand{\C}{\mathbb{C}}
\newcommand{\F}{\mathbb{F}}
\newcommand{\A}{\mathbb{A}}
\newcommand{\PP}{\mathbb{P}}
\newcommand{\E}{\mathbb{E}}
\newcommand{\eps}{\varepsilon}
\newcommand{\dd}{\,\mathrm{d}}
\begin{document}
\section*{UnsolvedMath ID: 30006291}
\textbf{Source problem code:} catalogue entry
\subsection{Definitions and mathematical statement}
A \emph{weak solution} is a solution interpreted through a weak formulation of the chosen fluid-mechanics PDE. An \emph{admissible weak solution} additionally satisfies the admissibility conditions imposed for that equation. The supplied record does not specify the PDE, weak formulation, function spaces, admissibility conditions, spatial domain, time interval, or initial and boundary data; consequently these terms cannot be defined more concretely from the supplied evidence.

\textbf{Research question.} Find a satisfactory solution concept for PDEs arising in mathematical fluid mechanics that addresses the nonuniqueness of admissible weak solutions, by resolving that nonuniqueness or selecting among such solutions.

The source presents this as a general research question. It supplies neither criteria defining ``satisfactory'' nor a precise requirement concerning existence, uniqueness, selection, stability, or agreement with classical solutions. It also does not assert nonuniqueness for every fluid-mechanics PDE or every choice of data. Thus a fully quantified mathematical formulation requires additional specifications; no particular equation, admissibility criterion, or selection rule is imposed in this draft.
\subsection{Short English statement}
The problem asks for a satisfactory notion of solution for fluid-mechanics PDEs in situations where admissible weak solutions are nonunique. The equations, data, admissibility conditions, and criteria for a successful solution concept remain unspecified, so the supplied question does not determine a single precise theorem.
\subsection{Sources}
\begin{itemize}
\item Supplied local registry record: \texttt{data/unsolvedmath/problems.json}, permanent UnsolvedMath ID \(30006291\), exact source problem code \texttt{catalogue entry}, title ``Selecting Among Nonunique Fluid-Mechanics Solutions.'' The supplied status is \texttt{open}; the supplied statement-status label is \texttt{corrected\_verified}. These are registry labels, not independent verification by this draft.
\item Supplied report citation: \emph{Mathematical Advances in Geophysical Fluid Dynamics} (2025), Oberwolfach Reports, pp.\ 1215--1270, DOI \texttt{10.4171/owr/2025/24}. Source URL and supplied literature-source link, labeled ``original OWR report'': \url{https://doi.org/10.4171/owr/2025/24}. The supplied background describes Oberwolfach Reports as records of mathematical workshops and their open-problem discussions.
\item Attribution in the supplied original statement: Simon Markfelder, joint work with Valentin Pellhammer, Christian Klingenberg, and Emil Wiedemann. The registry gives the proposed year as 2025 and does not separately identify a proposer.
\item Literature assessment dated \texttt{2026-08-22}: the supplied curation labels the question open and states that nonuniqueness of admissible weak solutions has motivated dissipative and selection frameworks, while the general search for a satisfactory solution concept remains unresolved. This is supplied dated curation, not independently verified; its time-sensitive status and cited claims require independent checking.
\end{itemize}
\end{document}
